% ============================================================
%  chapters/ch04-canonicalization.tex
% ============================================================

\chapter{Canonicalization and Equivalence}

\chapterepigraph{%
  Two paths that arrive at the same place\\
  are equivalent as routes\\
  but not as experiences.
}{Flyxion, \textit{Semantic Infrastructure}}

\sidenote{See App.\,A §A.4\\on equivalence\\relations; App.\,F\\on gesture\\canonicalization}

The previous three chapters established the foundational primitives:
accessibility (Chapter~1), containment (Chapter~2), and collapse
(Chapter~3). This chapter addresses a subtler question: when are two
different expressions, gestures, or trajectories the \emph{same thing}?

Sameness is not as simple as it appears. Two arithmetic expressions may
look different yet compute the same value. Two hand trajectories may take
different paths yet produce the same letter. Two scientific theories may
use different formalisms yet make the same predictions. In each case, a
\emph{canonicalization} procedure identifies the equivalence class to which
an object belongs — the unique representative of all the ways of being
the same thing.

% ── 4.1 ──────────────────────────────────────────────────────
\section{Why Different Trajectories Represent the Same Object}

Consider the number $12$. It can be represented as:
\[
  12, \quad 3 \times 4, \quad 2^2 \times 3, \quad 144/12,
  \quad 13 - 1, \quad \sum_{k=1}^{3} k \cdot (5-k), \quad \ldots
\]
These are not twelve different numbers. They are twelve different
\emph{trajectories} to the same number. The number $12$ is the equivalence
class of all arithmetic expressions that evaluate to $12$.

This observation, once made, reveals that the conventional picture of
mathematics as a study of objects is misleading. Mathematics is primarily
a study of \emph{equivalence classes of trajectories}. Numbers are not
primitive objects; they are the fixed points of evaluation processes.

\begin{definition}{Semantic Equivalence}{semantic-equiv}
  Two expressions $e_1, e_2 \in \mathcal{E}$ are \emph{semantically
  equivalent}, written $e_1 \equiv e_2$, if
  $\mathrm{val}(e_1) = \mathrm{val}(e_2)$.
  The \emph{equivalence class} of $e$ under $\equiv$ is
  $[e]_\equiv = \{e' \in \mathcal{E} : e' \equiv e\}$.
\end{definition}

The key insight is that semantics — the mathematical concept of
\emph{what something means} — is the study of these equivalence classes.
Two expressions have the same meaning precisely when they belong to
the same equivalence class. Meaning is the quotient of syntax by
equivalence.

% ── 4.2 ──────────────────────────────────────────────────────
\section{Gesture Equivalence}

The situation in gesture recognition is richer and more subtle than in
arithmetic, because gesture equivalence is not a sharp mathematical
relation — it is a \emph{graded} one.

\begin{definition}{Motor Equivalence Class}{motor-equiv}
  Two gesture trajectories $\gamma_1, \gamma_2 \in \gesturesp$ are
  \emph{motor-equivalent}, written $\gamma_1 \sim_m \gamma_2$, if
  they would be recognized as the same gesture by a competent observer
  under normal conditions. The \emph{motor equivalence class}
  $[\gamma]_m$ is the set of all trajectories equivalent to $\gamma$.
\end{definition}

The condition ``by a competent observer under normal conditions'' hides
considerable complexity. What counts as equivalent varies by context,
by the observer's expertise, by the ambient conditions (noise,
urgency, modality). Gesture equivalence is therefore not a fixed
mathematical relation but a \emph{family} of relations parametrized
by context.

\begin{remark}
  This contextual dependence is not a defect in the theory — it is
  the theory's most important feature. The same physical motion can
  be classified as the letter ``A'' in English handwriting, as the
  number ``4'' in another script, as a greeting gesture in yet another
  context. The gesture does not have an intrinsic meaning independent
  of context; it has meaning relative to a classification system.
  This is precisely the fiber-bundle structure: the same base point
  (physical motion) sits in different fibers depending on which
  projection (classification system) is applied.
\end{remark}

% ── 4.3 ──────────────────────────────────────────────────────
\section{Typing Equivalence}

Typing provides a particularly clean example of motor equivalence, because
the output space is discrete: every sequence of motor acts either produces
a valid character or it does not. This discreteness makes the equivalence
classes sharply defined.

\begin{example}
  The character ``a'' on a standard keyboard can be produced by:
  \begin{itemize}
    \item pressing the \texttt{A} key with the left pinky finger,
    \item pressing it with the left ring finger (awkward but possible),
    \item using an autocorrect expansion from ``aa'' to ``a'',
    \item pasting from clipboard containing ``a'',
    \item using a macro that types ``a''.
  \end{itemize}
  All five trajectories produce the same output character. They are
  motor-equivalent in the sense of producing identical symbolic output,
  but they are \emph{kinematically} very different: different fingers,
  different hand positions, different motor programs, different cognitive
  loads. The symbolic output ``a'' is the quotient that throws away
  all this kinematic variation.
\end{example}

The Gesture Before Symbol thesis (Chapter~10) argues that this throwing-away
is problematic. The kinematic variation that the discrete output discards
contains real information — about the typist's skill, emotional state,
attention, familiarity with the text, even identity (keystroke dynamics
are a biometric). Recognizing the equivalence class as a quotient makes
this information loss explicit and recoverable in principle.

% ── 4.4 ──────────────────────────────────────────────────────
\section{Mathematical Equivalence Classes}

The notion of equivalence class is one of the most powerful tools in
mathematics precisely because it makes the concept of ``sameness'' rigorous.

\begin{definition}{Quotient Set}{quotient-set}
  Given a set $X$ and an equivalence relation $\sim$ on $X$, the
  \emph{quotient set} $X/{\sim}$ is the set of equivalence classes:
  \[
    X/{\sim} \;=\; \bigl\{ [x]_\sim \,:\, x \in X \bigr\}
  \]
  The \emph{projection} $\pi : X \to X/{\sim}$ sends $x \mapsto [x]_\sim$.
\end{definition}

The quotient construction appears throughout mathematics whenever we want
to identify objects that ``look different but are the same.'' A few examples:

\begin{itemize}
  \item \textbf{Rational numbers:} $\mathbb{Q} = \mathbb{Z}^2 / {\sim}$
        where $(p,q) \sim (p',q')$ iff $pq' = p'q$. The fraction
        $\tfrac{2}{3}$ and $\tfrac{4}{6}$ are the same rational number
        — they are in the same equivalence class.
  \item \textbf{Topological spaces up to homeomorphism:} two topological
        spaces are equivalent if there is a continuous bijection between
        them with continuous inverse. The equivalence classes are
        topological types: a coffee mug and a donut are in the same
        class (both have genus 1).
  \item \textbf{Programs up to observational equivalence:} two programs
        are equivalent if they produce the same output for every input.
        This is the appropriate equivalence for program optimization.
  \item \textbf{Scientific theories up to empirical equivalence:}
        two theories are equivalent if they make the same predictions
        for all observable quantities. Underdetermination of theory
        by evidence is the non-triviality of this equivalence relation.
\end{itemize}

In each case, the quotient set captures the \emph{intrinsic} structure
while discarding irrelevant \emph{representational} variation.

% ── 4.5 ──────────────────────────────────────────────────────
\section{Canonical Forms}

For practical computation, it is useful to pick a single distinguished
representative from each equivalence class. This representative is the
\emph{canonical form}.

\begin{definition}{Canonical Form}{canonical-form}
  A \emph{canonicalization function} for an equivalence relation $\sim$
  on $X$ is a function $c : X \to X$ satisfying:
  \begin{enumerate}
    \item \textbf{Idempotence:} $c(c(x)) = c(x)$ for all $x \in X$.
    \item \textbf{Equivalence-preservation:}
          $x \sim y \Leftrightarrow c(x) = c(y)$.
  \end{enumerate}
  The \emph{canonical form} of $x$ is $c(x)$.
\end{definition}

The idempotence condition says: applying canonicalization twice is the same
as applying it once. Once something is canonical, it stays canonical. This
is the mathematical version of ``already in simplest form.''

\begin{example}
  For arithmetic expressions over integers:
  $c(3 \times 4) = c(2^2 \times 3) = c(12) = 12$.
  The canonical form is the fully evaluated integer. Evaluation is
  canonicalization.
\end{example}

\begin{example}
  For fractions: $c(p/q) = p'/q'$ where $\gcd(p',q') = 1$ and $q' > 0$.
  The canonical form is the reduced fraction in lowest terms.
\end{example}

\begin{example}
  For polynomials: $c(p)$ = coefficients written in decreasing degree
  with zero coefficients suppressed. The canonical form is the standard
  polynomial representation.
\end{example}

\begin{theorem}{Canonical Forms Classify}{canonical-classify}
  If $c : X \to X$ is a canonicalization function for $\sim$, then
  $X/{\sim} \cong c(X)$ (as sets). That is, the equivalence classes
  are in bijection with the canonical forms.
\end{theorem}

\begin{proof}
  Define $\phi : X/{\sim} \to c(X)$ by $\phi([x]) = c(x)$.
  This is well-defined: if $x \sim y$ then $c(x) = c(y)$ by
  equivalence-preservation. It is injective: if $c(x) = c(y)$
  then $x \sim y$ (by the same property), so $[x] = [y]$.
  It is surjective: for any $z \in c(X)$, $z = c(x)$ for some $x$,
  and $\phi([x]) = c(x) = c(c(x)) = c(z) = z$ by idempotence.
\end{proof}

% ── 4.6 ──────────────────────────────────────────────────────
\section{Fiber and Projection Viewpoints}

The quotient construction has a dual interpretation in terms of fibers
and projections, which is often geometrically illuminating.

\begin{definition}{Fiber of a Canonical Form}{fiber-canonical}
  Given a canonicalization $c : X \to X$, the \emph{fiber} over a
  canonical form $z \in c(X)$ is
  \[
    c^{-1}(z) = \{x \in X : c(x) = z\}
  \]
  — the set of all expressions that canonicalize to $z$.
\end{definition}

The fiber $c^{-1}(12)$ contains all arithmetic expressions that evaluate
to 12. The fiber $c^{-1}(\text{"a"})$ contains all motor sequences that
produce the character ``a''. The fiber is the \emph{extension} of the
canonical form: all the ways of being that thing.

\begin{figure}[h]
  \centering
  \begin{tikzpicture}[scale=1.0]
    % ============================================================
%  diagrams/fiber-projection-2d.tex
%  Fiber and projection: expression space onto canonical forms
% ============================================================

\tikzset{
  fiber/.style={RSVPmid!60, thin, opacity=0.7},
  projlabel/.style={font=\footnotesize\itshape, RSVPmid},
  canlabel/.style={font=\footnotesize\bfseries, RSVPdeep},
}

  % ── Top layer: expression space X ─────────────────────────
  \draw[RSVPdeep, thick, rounded corners=8pt]
    (-4, 0.4) rectangle (4, 2.8);
  \node[RSVPdeep, font=\small\scshape] at (-3.2, 2.5) {$X$};

  % Expressions (scattered points)
  \foreach \x/\y/\lbl in {
    -3.0/1.8/{$3\times4$},
    -1.8/2.2/{$2^2\!\times\!3$},
    -0.5/1.6/{$13-1$},
     0.5/2.3/{$\frac{36}{3}$},
     1.6/1.7/{$144/12$},
     3.0/2.1/{$\ldots$}
  }{
    \fill[RSVPaccent] (\x,\y) circle (2.5pt);
    \node[font=\tiny, above=1pt] at (\x,\y) {\lbl};
  }

  % A separate cluster for "7"
  \foreach \x/\y/\lbl in {
    -3.2/0.8/{$3+4$},
    -2.0/1.0/{$8-1$},
    -1.2/0.7/{$7^1$}
  }{
    \fill[RSVPpurple] (\x,\y) circle (2.5pt);
    \node[font=\tiny, above=1pt] at (\x,\y) {\lbl};
  }

  % A cluster for "5"
  \foreach \x/\y/\lbl in {
     2.0/0.8/{$2+3$},
     2.8/1.1/{$10/2$},
     3.2/0.6/{$5^1$}
  }{
    \fill[RSVPmid] (\x,\y) circle (2.5pt);
    \node[font=\tiny, above=1pt] at (\x,\y) {\lbl};
  }

  % ── Projection arrow ──────────────────────────────────────
  \draw[-{Stealth[length=8pt]}, RSVPdeep, very thick]
    (0, 0.3) -- (0, -0.7)
    node[midway, right=4pt, font=\small\itshape, RSVPdeep] {$c$};

  % ── Bottom layer: canonical forms c(X) ────────────────────
  \draw[RSVPdeep, thick, rounded corners=8pt]
    (-4, -1.8) rectangle (4, -0.9);
  \node[RSVPdeep, font=\small\scshape] at (-3.2, -1.0) {$c(X)$};

  % Canonical points
  \fill[RSVPpurple] (-2.5,-1.35) circle (3.5pt);
  \node[canlabel, below=3pt] at (-2.5,-1.35) {$7$};

  \fill[RSVPaccent] (0,-1.35) circle (3.5pt);
  \node[canlabel, below=3pt] at (0,-1.35) {$12$};

  \fill[RSVPmid] (2.5,-1.35) circle (3.5pt);
  \node[canlabel, below=3pt] at (2.5,-1.35) {$5$};

  % ── Fiber lines ───────────────────────────────────────────
  % Fibers for 12 (amber cluster)
  \foreach \x/\y in {-3.0/1.8, -1.8/2.2, -0.5/1.6, 0.5/2.3, 1.6/1.7, 3.0/2.1}{
    \draw[fiber, RSVPaccent!70]
      (\x,\y) .. controls (\x*0.5, 0.0) and (0, -0.5) .. (0,-0.9);
  }
  % Fibers for 7 (purple)
  \foreach \x/\y in {-3.2/0.8, -2.0/1.0, -1.2/0.7}{
    \draw[fiber, RSVPpurple!70]
      (\x,\y) .. controls (\x*0.7, 0.0) and (-2.5,-0.6) .. (-2.5,-0.9);
  }
  % Fibers for 5 (teal)
  \foreach \x/\y in {2.0/0.8, 2.8/1.1, 3.2/0.6}{
    \draw[fiber, RSVPmid!70]
      (\x,\y) .. controls (\x*0.7, 0.0) and (2.5,-0.6) .. (2.5,-0.9);
  }

  % ── Fiber label ───────────────────────────────────────────
  \draw[RSVPaccent!50, thin, dashed]
    (-3.2,0.35) -- (3.2,0.35) -- (3.2,2.85) -- (-3.2,2.85) -- cycle;
  \node[RSVPaccent!80, font=\footnotesize\itshape] at (4.6, 1.6)
    {$c^{-1}(12)$};
  \draw[RSVPaccent!50, thin] (4.1,1.6) -- (3.3,1.6);

  \end{tikzpicture}
  \caption{Fiber and projection geometry. The expression space $X$ (top)
    contains many expressions that canonicalize to the same value in $c(X)$
    (bottom). The projection $c$ collapses each fiber to a point.
    The fiber over the canonical form $z = 12$ contains all expressions
    that evaluate to $12$: $(3 \times 4)$, $(2^2 \times 3)$, $(13-1)$,
    and infinitely many others.}
  \label{fig:fiber-projection}
\end{figure}

The fiber viewpoint makes the information-loss interpretation of
canonicalization explicit (cf.\  Appendix~B). The canonical form retains
only the equivalence class membership; the fiber contains all the
additional information that was discarded.

% ── 4.7 ──────────────────────────────────────────────────────
\section{Canonicalization in the RSVP Framework}

Within RSVP, canonicalization has a field-theoretic interpretation. The
accessibility entropy $S$ measures the volume of admissible futures.
A canonical form is a state of locally \emph{minimal} entropy within
its equivalence class — the state that has already collapsed as much
variation as possible.

\begin{definition}{RSVP Canonical State}{rsvp-canonical}
  A state $x \in X$ is \emph{RSVP-canonical} within its equivalence
  class $[x]_\sim$ if it minimizes the constraint violation functional
  $\kappa$ and is a local minimum of accessible future volume
  variation:
  \[
    x_\text{can} = \arg\min_{y \sim x} \kappa(y).
  \]
\end{definition}

Under accessibility relaxation dynamics (Appendix~J), all trajectories
within an equivalence class converge to the canonical representative.
The dynamics enforce canonicalization automatically: given enough time,
equivalent states converge to their shared canonical form.

\begin{remark}
  This is why stable structures persist. A galaxy, an organism, a crystal,
  a mathematical concept — all are RSVP-canonical states: local minima
  of constraint violation within the equivalence class of states with
  similar accessibility properties. They persist not because they are
  special objects, but because they are the attractors of the
  canonicalization dynamics.
\end{remark}

% ── 4.8 ──────────────────────────────────────────────────────
\section{The Canonicalization Principle}

The various forms of canonicalization discussed in this chapter — arithmetic
evaluation, gesture recognition, typing, mathematical reduction, RSVP
relaxation — are all instances of a single principle:

\begin{namedthm}{The Canonicalization Principle}
  Any system organized by equivalence classes has a natural dynamics:
  the motion of equivalent states toward their shared canonical
  representative. Stable structures are the canonical representatives
  of their equivalence class under the ambient dynamics.
\end{namedthm}

This principle connects the first four chapters into a single coherent
picture. Objects (Chapter~1) are equivalence classes of accessible
trajectories. Boundaries (Chapter~2) define the admissible region within
which canonical forms reside. Collapse (Chapter~3) is the mechanism by
which equivalence classes are reduced to their representatives.
Canonicalization (this chapter) is the principle that unifies the three:
the persistent structure that remains when collapse is complete.

% ── Exercises ────────────────────────────────────────────────
\section*{Exercises}

\begin{exercise}
  Let $X = \{$all finite arithmetic expressions over $\mathbb{Z}\}$
  with $e_1 \equiv e_2$ iff they evaluate to the same integer.
  Show that $\equiv$ is an equivalence relation. Describe the
  canonicalization function $c$ and verify that it satisfies
  idempotence and equivalence-preservation.
\end{exercise}

\begin{exercise}
  The rational numbers $\mathbb{Q}$ are defined as the quotient
  $(\mathbb{Z} \times \mathbb{Z}^*)/{\sim}$ where $(p,q) \sim (p',q')$
  iff $pq' = p'q$. Verify that $\sim$ is an equivalence relation.
  Define the canonical form for fractions (reduced form). Show that
  every equivalence class has exactly one canonical representative.
\end{exercise}

\begin{exercise}
  In the RSVP framework, accessibility entropy $S(x) = \log \Omega(x)$
  measures future volume. Argue (informally) why two states in the same
  equivalence class — making the same predictions for all observables —
  should have the same $S$ value. Under what circumstances might they
  have different $S$ values?
\end{exercise}

\begin{exercise}[title={Open problem}]
  Define a notion of \emph{canonical gesture} for handwriting that
  is consistent with the motor equivalence relation $\sim_m$. What
  properties should the canonical representative of $[\gamma]_m$
  have? (Consider: minimum energy, minimum duration, maximum robustness
  to perturbation.) Are these criteria consistent with each other?
  If not, which should take precedence and why?
\end{exercise}
